% GENERATED FILE. Edit canonical lessons, then run npm run generate:books.
% canonical-source-hash: fnv1a64:19749d0c758e1e31
% canonical-lessons: FA-C446-jahat, FA-C446-qotbnama, FA-C446-kilometr, FA-C446-masafat, FA-C446-dur-shodan

\chapter{Directions and distances}
\label{ch:fa-directions-and-distances}
\hlchaptermodality{446}

\begin{chapteropening}
\textbf{By the end of this chapter:} \emph{I can talk about direction and distance: a direction, a compass, kilometres, distance and moving away.}

Complete the last lesson of chapter 446: I can talk about direction and distance: a direction, a compass, kilometres, distance and moving away.
\end{chapteropening}

\section[\texorpdfstring{\fa{جهت} (jahat)}{جهت (jahat)}]{\fa{جهت} (jahat) — a direction; a side}

\label{lesson:FA-C446-jahat}



\begin{quote}
Before the new one: say the Persian for a fork, then the Persian for to cross.
\end{quote}

\subsection*{You'll want to know: \fa{جهت}}
\textbf{\fa{جهت}} (\emph{jahat}) — \textquotedblleft{}a direction; a side\textquotedblright{}.

\begin{quote}
\fa{ما} \fa{در} \fa{جهت} \fa{اشتباه} \fa{رانندگی} \fa{میکردیم}.
\end{quote}

\begin{quote}
\emph{Mâ dar jahat-e eshtebâh rânandegi mikardim.}
\end{quote}

\begin{quote}
We were driving in the wrong direction.
\end{quote}

Arabic plural \emph{jahât}. \emph{Be in jahat} — for this reason.

\begin{cousinweb}
\emph{jahat} was borrowed from Arabic and is usually traced to Arabic \emph{jiha} 'side, direction'. EtymDB does not list it.
\end{cousinweb}

\begin{grammarlens}[title={What you've built}]
One more word for telling, explaining and describing at B1.
\end{grammarlens}

\subsection*{Your turn}
\begin{itemize}
\raggedright
  \item \emph{Say it:} \emph{jahat}
  \item \emph{Say it:} \emph{jahat}, once more
  \item \emph{Say it:} a sentence of your own with \emph{jahat}
  \item \emph{Recall:} say the Persian for a fork, then the Persian for to cross, then say \emph{jahat} again
\end{itemize}

\begin{tcolorbox}[breakable,colback=teal!4,colframe=teal!35!black,title={Before you move on}]
What does \fa{جهت} mean? (\textquotedblleft{}A direction; a side\textquotedblright{}.) Say it once more.
\end{tcolorbox}
\section[\texorpdfstring{\fa{قطب} \fa{نما} (qotbnamâ)}{قطب نما (qotbnamâ)}]{\fa{قطب} \fa{نما} (qotbnamâ) — a compass}

\label{lesson:FA-C446-qotbnama}



\begin{quote}
Before the new one: say the Persian for to cross, then the Persian for a direction.
\end{quote}

\subsection*{You'll want to know: \fa{قطب} \fa{نما}}
\textbf{\fa{قطب} \fa{نما}} (\emph{qotbnamâ}) — \textquotedblleft{}a compass\textquotedblright{}.

\begin{quote}
\fa{در} \fa{کوه} \fa{با} \fa{قطب} \fa{نما} \fa{راه} \fa{را} \fa{پیدا} \fa{کردیم}.
\end{quote}

\begin{quote}
\emph{Dar kuh bâ qotbnamâ râh râ peydâ kardim.}
\end{quote}

\begin{quote}
In the mountains we found the way with a compass.
\end{quote}

Literally 'pole-shower': \emph{qotb} 'pole' + \emph{namâ} 'showing'.

\begin{cousinweb}
\emph{qotbnamâ} is built from \emph{qotb} 'pole', borrowed from Arabic \emph{quṭb} (EtymDB lists the Arabic word), and native \emph{namâ} 'showing', the stem of \emph{nemudan} 'to show'.
\end{cousinweb}

\begin{grammarlens}[title={What you've built}]
One more word for telling, explaining and describing at B1.
\end{grammarlens}

\subsection*{Your turn}
\begin{itemize}
\raggedright
  \item \emph{Say it:} \emph{qotbnamâ}
  \item \emph{Say it:} \emph{qotbnamâ}, once more
  \item \emph{Say it:} a sentence of your own with \emph{qotbnamâ}
  \item \emph{Recall:} say the Persian for to cross, then the Persian for a direction, then say \emph{qotbnamâ} again
\end{itemize}

\begin{tcolorbox}[breakable,colback=teal!4,colframe=teal!35!black,title={Before you move on}]
What does \fa{قطب} \fa{نما} mean? (\textquotedblleft{}A compass\textquotedblright{}.) Say it once more.
\end{tcolorbox}
\section[\texorpdfstring{\fa{کیلومتر} (kilometr)}{کیلومتر (kilometr)}]{\fa{کیلومتر} (kilometr) — a kilometre}

\label{lesson:FA-C446-kilometr}



\begin{quote}
Before the new one: say the Persian for a direction, then the Persian for a compass.
\end{quote}

\subsection*{You'll want to know: \fa{کیلومتر}}
\textbf{\fa{کیلومتر}} (\emph{kilometr}) — \textquotedblleft{}a kilometre\textquotedblright{}.

\begin{quote}
\fa{تا} \fa{شهر} \fa{بعدی} \fa{صد} \fa{کیلومتر} \fa{راه} \fa{است}.
\end{quote}

\begin{quote}
\emph{Tâ shahr-e ba'di sad kilometr râh ast.}
\end{quote}

\begin{quote}
It is a hundred kilometres to the next town.
\end{quote}

Speed is \emph{kilometr dar sâ'at}.

\begin{cousinweb}
\emph{kilometr} was borrowed from French \emph{kilomètre} (EtymDB), which was built from Greek \emph{khílioi} 'a thousand' and \emph{métron} 'measure'.
\end{cousinweb}

\begin{grammarlens}[title={What you've built}]
One more word for telling, explaining and describing at B1.
\end{grammarlens}

\subsection*{Your turn}
\begin{itemize}
\raggedright
  \item \emph{Say it:} \emph{kilometr}
  \item \emph{Say it:} \emph{kilometr}, once more
  \item \emph{Say it:} a sentence of your own with \emph{kilometr}
  \item \emph{Recall:} say the Persian for a direction, then the Persian for a compass, then say \emph{kilometr} again
\end{itemize}

\begin{tcolorbox}[breakable,colback=teal!4,colframe=teal!35!black,title={Before you move on}]
What does \fa{کیلومتر} mean? (\textquotedblleft{}A kilometre\textquotedblright{}.) Say it once more.
\end{tcolorbox}
\section[\texorpdfstring{\fa{مسافت} (masâfat)}{مسافت (masâfat)}]{\fa{مسافت} (masâfat) — a distance (to be travelled)}

\label{lesson:FA-C446-masafat}



\begin{quote}
Before the new one: say the Persian for a compass, then the Persian for a kilometre.
\end{quote}

\subsection*{You'll want to know: \fa{مسافت}}
\textbf{\fa{مسافت}} (\emph{masâfat}) — \textquotedblleft{}a distance (to be travelled)\textquotedblright{}.

\begin{quote}
\fa{مسافت} \fa{خانه} \fa{تا} \fa{مدرسه} \fa{زیاد} \fa{نیست}.
\end{quote}

\begin{quote}
\emph{Masâfat-e khâne tâ madrese ziyâd nist.}
\end{quote}

\begin{quote}
The distance from home to school isn't great.
\end{quote}

Close to \emph{fâsele}, but used for a journey's length.

\begin{cousinweb}
\emph{masâfat} was borrowed from Arabic and is usually traced to Arabic \emph{masāfa} 'distance, the length of a journey'. EtymDB does not list it.
\end{cousinweb}

\begin{grammarlens}[title={What you've built}]
One more word for telling, explaining and describing at B1.
\end{grammarlens}

\subsection*{Your turn}
\begin{itemize}
\raggedright
  \item \emph{Say it:} \emph{masâfat}
  \item \emph{Say it:} \emph{masâfat}, once more
  \item \emph{Say it:} a sentence of your own with \emph{masâfat}
  \item \emph{Recall:} say the Persian for a compass, then the Persian for a kilometre, then say \emph{masâfat} again
\end{itemize}

\begin{tcolorbox}[breakable,colback=teal!4,colframe=teal!35!black,title={Before you move on}]
What does \fa{مسافت} mean? (\textquotedblleft{}A distance (to be travelled)\textquotedblright{}.) Say it once more.
\end{tcolorbox}
\section[\texorpdfstring{\fa{دور} \fa{شدن} (dur shodan)}{دور شدن (dur shodan)}]{\fa{دور} \fa{شدن} (dur shodan) — to move away, to get further away}

\label{lesson:FA-C446-dur-shodan}



\begin{quote}
Before the new one: say the Persian for a kilometre, then the Persian for a distance.
\end{quote}

\subsection*{You'll want to know: \fa{دور} \fa{شدن}}
\textbf{\fa{دور} \fa{شدن}} (\emph{dur shodan}) — \textquotedblleft{}to move away, to get further away\textquotedblright{}.

\begin{quote}
\fa{قایق} \fa{کم} \fa{کم} \fa{از} \fa{ساحل} \fa{دور} \fa{شد}.
\end{quote}

\begin{quote}
\emph{Qâyeq kam-kam az sâhel dur shod.}
\end{quote}

\begin{quote}
The boat slowly moved away from the shore.
\end{quote}

\emph{dur} 'far' + \emph{shodan}. The opposite is \emph{nazdik shodan}.

\begin{cousinweb}
\emph{dur shodan} is built from native \emph{dur} 'far', which EtymDB traces to Proto-Iranian \emph{duHráh} 'far' and PIE \emph{duh\textsubscript{2}-ró-s} 'long', and the native verb \emph{shodan} 'to become'.
\end{cousinweb}

\begin{grammarlens}[title={What you've built}]
One more word for telling, explaining and describing at B1.
\end{grammarlens}

\subsection*{Your turn}
\begin{itemize}
\raggedright
  \item \emph{Say it:} \emph{dur shodan}
  \item \emph{Say it:} \emph{dur shodan}, once more
  \item \emph{Say it:} a sentence of your own with \emph{dur shodan}
  \item \emph{Recall:} say the Persian for a kilometre, then the Persian for a distance, then say \emph{dur shodan} again
\end{itemize}

\begin{tcolorbox}[breakable,colback=teal!4,colframe=teal!35!black,title={Before you move on}]
What does \fa{دور} \fa{شدن} mean? (\textquotedblleft{}To move away, to get further away\textquotedblright{}.) Say it once more.
\end{tcolorbox}
